\chapter{Groin Irritation and Infection}
\index{Groin Irritation and Infection}

\section*{Definition and Overview}
Groin irritation and infection refers to a range of common skin conditions affecting the folds of skin in the groin area, where warmth, moisture, and friction make the skin vulnerable. The most common causes include fungal infection (tinea cruris, commonly called jock itch), bacterial infections, irritation from friction or personal products, and inflammatory skin conditions such as eczema and psoriasis. Symptoms include itching, redness, a rash, burning, and sometimes soreness or small bumps. Most causes of groin irritation are easily treated with good hygiene, keeping the area dry, and appropriate creams or medicines, and most settle within a few weeks.

\section*{Epidemiology}
Groin skin problems are very common. Tinea cruris (jock itch) is one of the most frequent, particularly in men, in athletes, and in people who sweat heavily or wear tight clothing. Around 10--20\% of adults will experience a groin fungal infection at some point. Bacterial skin infections, including boils and folliculitis, also commonly affect the groin, and intertrigo (chafing in the skin folds) is common in hot weather, in people who are overweight, and in people with diabetes. Groin irritation is more common in men than women.

\section*{Aetiology and Pathophysiology}
The groin is prone to skin problems because of its anatomy and environment.

\begin{itemize}
    \item \textbf{Moisture and warmth:} the skin folds trap sweat and heat, favouring fungal and bacterial growth
    \item \textbf{Friction:} rubbing of skin against skin or clothing damages the skin barrier
    \item \textbf{Fungal infection:} dermatophyte fungi cause tinea cruris, which can spread from the feet (tinea pedis)
    \item \textbf{Bacterial infection:} \textit{Staphylococcus aureus} and other bacteria cause folliculitis, boils, and cellulitis
    \item \textbf{Skin conditions:} eczema, psoriasis, and contact dermatitis from soaps, detergents, or deodorants
    \item \textbf{Contributing factors:} obesity, diabetes, poor hygiene, tight clothing, and prolonged sitting
    \item \textbf{Spread:} scratching can spread infection and introduce bacteria
\end{itemize}

\section*{Clinical Features}
Symptoms vary with the cause, and the pattern of the rash is often diagnostic.

\begin{itemize}
    \item Itching, which can be intense, especially in fungal infection
    \item A red rash in the skin folds, often with a distinct edge in tinea cruris
    \item Burning or stinging, especially after exercise or when the skin is wet
    \item Soreness and cracked skin
    \item Small pus-filled bumps (folliculitis) or boils
    \item Flaking, scaling, or oozing skin
    \item An unpleasant odour in bacterial overgrowth
    \item Redness and swelling that spreads beyond the fold in cellulitis
\end{itemize}

\section*{Differential Diagnosis}
Because several conditions look similar, accurate diagnosis guides treatment.

\begin{itemize}
    \item Tinea cruris (jock itch): a fungal infection with a well-defined, scaly edge
    \item Candidal intertrigo: a yeast infection that appears as a moist, red rash with satellite spots
    \item Intertrigo: simple chafing from moisture and friction without infection
    \item Contact dermatitis from soaps, detergents, or fabric
    \item Eczema and psoriasis affecting the groin
    \item Folliculitis and boils from bacterial infection
    \item Cellulitis, a deeper bacterial infection that spreads
    \item Hidradenitis suppurativa, a chronic condition causing painful lumps in the skin folds
    \item Sexually transmitted infections, which can cause groin symptoms
    \item Scabies and other infestations, which can affect the groin
\end{itemize}

\section*{When to Seek Medical Care}
See a doctor if the rash is severe, painful, spreading, or not improving after a week of self-care, or if you have diabetes or a weakened immune system.

\textbf{Contact your local emergency services for:}
\begin{itemize}
    \item Spreading redness, warmth, pain, or swelling with fever, which may indicate cellulitis
    \item A painful, swollen lump in the groin with fever
    \item Difficulty passing urine or a scrotum that is suddenly very swollen and painful
    \item Signs of sepsis, including high fever, confusion, or collapse
\end{itemize}

\section*{Diagnosis and Investigations}
Most groin irritation is diagnosed clinically, with tests when the cause is unclear.

\begin{itemize}
    \item \textbf{Clinical examination:} the appearance and location of the rash often identify the cause
    \item \textbf{Skin scraping:} examined under a microscope for fungal elements
    \item \textbf{Culture:} swabs of weeping or pustular areas identify bacteria or yeast
    \item \textbf{Blood glucose testing:} recurrent or severe infections can be the first sign of diabetes
    \item \textbf{Allergy testing:} if contact dermatitis is suspected
    \item \textbf{Referral:} to a dermatologist for persistent or unusual cases
\end{itemize}

\section*{Management}
Treatment depends on the cause, and hygiene measures are important in all cases.

\begin{itemize}
    \item \textbf{Fungal infection:} antifungal creams such as clotrimazole or terbinafine, applied for at least two weeks
    \item \textbf{Yeast infection:} antifungal creams active against candida
    \item \textbf{Bacterial infection:} antiseptic washes and, if needed, antibiotic creams or oral antibiotics
    \item \textbf{Eczema and dermatitis:} moisturisers, mild steroid creams, and avoidance of irritants
    \item \textbf{Keeping the area dry:} washing daily, drying thoroughly, and wearing loose cotton underwear
    \item \textbf{Treating the feet:} tinea on the feet should be treated to prevent reinfection of the groin
    \item \textbf{Managing contributing factors:} weight management and blood glucose control
    \item \textbf{Severe infection:} oral antibiotics for cellulitis or boils
\end{itemize}

\section*{Complications}
\begin{itemize}
    \item Spread of fungal infection to the thighs, buttocks, or feet
    \item Secondary bacterial infection of scratched skin
    \item Cellulitis requiring antibiotics
    \item Abscess formation from boils
    \item Chronic or recurrent infection when underlying factors are not addressed
    \item Persistent itching affecting sleep and quality of life
    \item Hyperpigmentation (darkening of the skin) that can persist after healing
\end{itemize}

\section*{Prognosis and Outcomes}
Most causes of groin irritation and infection settle within one to four weeks with correct treatment. Fungal infections respond to antifungals but recur easily if the area is kept moist, so ongoing hygiene is important. Bacterial infections clear with antibiotics, and inflammatory conditions are controlled with appropriate creams. Chronic conditions such as hidradenitis suppurativa and diabetes-related infections need ongoing management, but most people achieve good control and resolution of symptoms.

\section*{Prevention and Self-Care}
\begin{itemize}
    \item Wash the groin daily with plain water or mild soap, and dry thoroughly
    \item Wear loose-fitting cotton underwear and change out of sweaty clothes promptly
    \item Avoid sharing towels, and dry the groin before the feet
    \item Use a separate towel for the groin and feet, or dry the feet last
    \item Treat athlete's foot promptly to prevent spread to the groin
    \item Avoid harsh soaps, deodorants, and fabric softeners in the groin area
    \item Apply antifungal powder to keep the area dry if you are prone to fungal infections
    \item Maintain a healthy weight and control diabetes
    \item See a doctor if the rash persists, worsens, or is associated with pain or fever
\end{itemize}

\section*{Sources and Further Reading}
The following reputable sources provide further information for healthcare students and patients seeking reliable, current advice about groin irritation and infection.

\begin{itemize}
    \item Healthdirect Australia. \textit{Jock itch (tinea cruris).} https://www.healthdirect.gov.au/
    \item DermNet NZ. \textit{Groin skin conditions.} https://dermnetnz.org/
    \item Australasian College of Dermatologists. \textit{Skin infections information.} https://www.dermcoll.edu.au/
    \item NSW Health. \textit{Fungal skin infections fact sheet.} https://www.health.nsw.gov.au/
    \item NHS. \textit{Jock itch and groin rashes.} https://www.nhs.uk/
    \item Mayo Clinic. \textit{Jock itch: symptoms and causes.} https://www.mayoclinic.org/
\end{itemize}
